\begin{claim}{D.77}{Proved}{A proper three-section subgroup with every primary image full}
There is an explicit choice of one representative of each
of the three section classes in the standard pants extension
for which the specified incidence map of D.76 is not onto,
although the chosen orientation-kernel generators have full
image in every maximal pro-$p$ quotient.

In the full profinite coproduct
\[
 E_0=C_2(a)\amalg C_2(b)\amalg C_2(c),\qquad
 F_0=\ker\bigl(E_0\longrightarrow C_2,\ a,b,c\longmapsto1\bigr),
                                                               \tag{D.77.1}
\]
put
\[
 x=ba,\quad y=ca,\quad
 \delta=\operatorname{Inn}(a)|_{F_0},\quad
 g=x[x,y]x^{-1},\quad
 Y=g\,y\,\delta(g)^{-1}.
                                                               \tag{D.77.2}
\]
Then $F_0=\widehat F(x,y)$, $\delta$ inverts $x,y$, and
\[
 \begin{gathered}
 T_0=\langle a\rangle,\qquad T_1=\langle b\rangle,\qquad
 T_2=\langle gcg^{-1}\rangle,\\
 H=\overline{\langle T_0,T_1,T_2\rangle}
       =K\rtimes\langle a\rangle,\qquad
 K=\overline{\langle x,Y\rangle}=H\cap F_0 .
 \end{gathered}                                                \tag{D.77.3}
\]
The pair $x,Y$ induces the identity on
$F_0^{\rm ab}=\widehat\Z^2$ and generates $F_0^{(p)}$
for every prime $p$, but $K$ is a proper subgroup of $F_0$
and $H$ is a proper subgroup of $E_0$.
The corresponding section augmentation module is the free
module of rank two in D.76. Its specified incidence cokernel
has a quotient of dimension five over $\mathbb F_2$.

The same finite quotient proving properness proves that
$Y$ is not conjugate to any full unit power of the original
$y$. Thus this example retains the three section labels
and all primary generation tests, and explicitly fails the
additional original-cusp alignment in D.76's distinguished
triple.
\end{claim}
\begin{proof}
\emph{The full extension and inverse maps.}
Let $F=\widehat F(x,y)$ and let $\delta$ be the actual
automorphism inverting its named generators.
The mutually inverse full profinite homomorphisms are
\[
 \begin{aligned}
 E_0&\longrightarrow F\rtimes_\delta C_2(a),
       &(a,b,c)&\longmapsto(a,xa,ya),\\
 F\rtimes_\delta C_2(a)&\longrightarrow E_0,
       &(x,y,a)&\longmapsto(ba,ca,a).
 \end{aligned}                                                \tag{D.77.4}
\]
The three images in the first row square to one.
The second row preserves
$a(ba)a=(ba)^{-1}$ and $a(ca)a=(ca)^{-1}$.
Both composites fix every indicated generator.
Their quotient characters to $C_2$ agree on those
generators, so their restriction identifies the actual
kernel in (D.77.1) with the full $F$, including its topology.

The subgroups $\langle a\rangle,\langle b\rangle,\langle c\rangle$
are the three section classes for this standard inversion
extension, as in D.34. Their images in the actual
abelian quotient
\[
 E_0\longrightarrow C_2^3,\qquad
       a\longmapsto e_0,\quad b\longmapsto e_1,\quad
       c\longmapsto e_2
                                                               \tag{D.77.5}
\]
are distinct and unchanged by conjugation.
Replacing $\langle c\rangle$ by its conjugate $T_2$
therefore retains exactly one representative of each class.

Since $c=ya$, direct multiplication gives
\[
 gcg^{-1}=gya g^{-1}
        =g y\delta(g)^{-1}a=Ya,\qquad
 \delta(Y)=Y^{-1}.
                                                               \tag{D.77.6}
\]
Thus $K$ is $\delta$-invariant. The identities
$x=ba$ and $Y=(gcg^{-1})a$ show that the group generated
by the three $T_i$ is exactly $K\langle a\rangle$.
The quotient label of $a$ is nonzero, whereas $K\leq F_0$;
hence the intersection and semidirect-product statements
in (D.77.3) follow.

\emph{The actual all-prime images.}
The word $g=x[x,y]x^{-1}$ lies in $F_0'$.
Taking full abelianizations in (D.77.2) gives
\[
                     x^{\rm ab}=(1,0),\qquad
                     Y^{\rm ab}=(0,1).
                                                               \tag{D.77.7}
\]
For each prime $p$, the continuous map
$x\mapsto x,\ y\mapsto Y$ induces an endomorphism
of the free pro-$p$ group $F_0^{(p)}$.
Its image maps onto the elementary abelian quotient
$F_0^{(p)}/\Phi(F_0^{(p)})$ by (D.77.7).
To verify full generation here, a proper closed image
would have proper image in some finite $p$-group quotient.
A maximal subgroup of that finite $p$-group has index $p$,
so a nonzero continuous character to $\mathbb F_p$
would vanish on both specified generators.
That contradicts (D.77.7).
The induced endomorphism is therefore onto.

It is also injective. For completeness, a finitely
generated pro-$p$ group has only finitely many
homomorphisms to the finitely many $p$-groups of any
bounded order. The intersections of their kernels
are open fully invariant subgroups and separate points.
An onto endomorphism induces a bijection on each
of these finite quotients. Its kernel is therefore
contained in their trivial intersection.
This proves the asserted automorphism on every
maximal pro-$p$ quotient by the actual specified maps.
Equation (D.77.6) also gives
$\delta\rho=\rho\delta$ for the full endomorphism
$\rho(x)=x,\ \rho(y)=Y$.

\emph{The finite quotient detecting both proper inclusions.}
Use permutations of $\{1,2,3,4,5\}$ with multiplication
$(uv)(i)=u(v(i))$. Define the full continuous quotient
$q:E_0\to A_5$ on the three order-two factors by
\[
 \begin{aligned}
 q(a)&=(1\,2)(3\,4),&
 q(b)&=(1\,3)(2\,5),&
 q(c)&=(1\,4)(2\,3).
 \end{aligned}                                                \tag{D.77.8}
\]
The coproduct universal property gives this homomorphism.
The exact finite multiplication and generation
calculation gives
\[
 \begin{gathered}
 q(x)=(1\,5\,2\,3\,4),\qquad
 q(y)=(1\,3)(2\,4),\qquad q(g)=(1\,3\,2),\\
 q(gcg^{-1})=q(a),\qquad q(Y)=1,\\
 \langle q(a),q(b),q(c)\rangle
       =\langle q(x),q(y)\rangle=A_5 .
 \end{gathered}                                                \tag{D.77.9}
\]
These equalities are also exhaustively verified by
the finite producer below: closure under multiplication
lists exactly sixty elements in each of the last two
generated groups.
In particular the restriction $q|_{F_0}$ is onto.

The images of the chosen subgroups are instead
\[
 q(H)=\langle q(a),q(b)\rangle=D_{10},\qquad
 q(K)=\langle q(x)\rangle=C_5 .
                                                               \tag{D.77.10}
\]
Indeed $q(x)=q(b)q(a)$ has order five and
$q(a)$ conjugates it to its inverse. The five rotations
and five reflections are distinct, proving the first
order in (D.77.10); the second is immediate.
Both are proper subgroups of the specified image $A_5$.
This proves $H\ne E_0$ and $K\ne F_0$ in the full
profinite groups.

\emph{The specified incidence cokernel and cusp comparison.}
Let $k=\mathbb F_2$, $R=k[[E_0]]$, and use the actual
$T_i$ from (D.77.3). D.76 applies to give
$\Delta_T\cong R^2$ and
\[
 \begin{aligned}
 d_T:R^2&\longrightarrow\Delta_T,\qquad
 (r_1,r_2)\longmapsto r_1(e_1-e_0)+r_2(e_2-e_0),\\
 \operatorname{coker}d_T
      &\cong\ker\bigl(k[[E_0/H]]\to k\bigr).
 \end{aligned}                                                \tag{D.77.11}
\]
There is a specified continuous surjection of coset spaces
\[
 E_0/H\longrightarrow A_5/D_{10},\qquad
                         zH\longmapsto q(z)D_{10}.
                                                               \tag{D.77.12}
\]
It is well defined by (D.77.10) and is onto by (D.77.9).
The induced permutation-module surjection respects
augmentation. Its restriction maps the kernel in
(D.77.11) onto
$\ker(k[A_5/D_{10}]\to k)$, which has dimension
$6-1=5$. This is an explicit nonzero finite quotient
of the actual incidence cokernel.

Finally $q(y)$ has order two, while $q(Y)=1$.
For every $\mu\in\widehat\Z^\times$, its residue modulo
two is one, so $q(y^\mu)=q(y)$ still has order two.
A conjugacy $Y=w y^\mu w^{-1}$ in $F_0$ would induce
a conjugacy between elements of orders one and two
under $q|_{F_0}$, which is impossible.
This proves the claimed failure of original-y alignment
through the same full homomorphism.
D.76's distinguished element $hyh^{-1}$ retains that
alignment by its defining equation; this example does
not solve its remaining generation or third-cusp problem.

The producer
\path{code/part_c/pants_incidence_limit.py}
executes twenty exact checks: the free-word expression
for $g$, its two zero abelian coordinates,
$\delta(Y)=Y^{-1}$, the identity abelianization
of $x,Y$, and all finite permutations, element orders
and generated image orders in (D.77.8)--(D.77.10).
The preceding Frattini argument supplies the
all-prime assertion; no finite list of primes is substituted.
\end{proof}
